\chapter*{Adnotare}
\author{Vasile Drumea}
\thispagestyle{empty}
\par Teza de masterat cu titlul "Sistem de monitorizare a transportului public în baza tehnologiei GPS", a fost elaborată de studentul grupei CRI-201M, Vasile Drumea, Universitatea Tehnică a Moldovei. Acest proiect este alcătuit din 3 capitole, introducere, concluzii și bibliografie.
\par \textbf{Cuvinte cheie:} IoT, MQTT, GPS, Transport, Comunicație, Dispozitiv, Publish-Subscribe.
\par Transportul public constituie un element esențial pentru majoritatea populației din întreaga lume. Indiferent de vârstă, ocupație sau stare materială, zilnic oamenii au nevoie de o modalitate de transport pe care se pot baza în orice moment, respectiv sistemul de transport public trebuie să fie organizat pentru a satisface cererea. De asemenea, utilizatorii au nevoie de surse de informare cât mai actuale și relevante, pentru a putea lua decizii cu privire la transport. Deseori pentru același traseu pot exista multiple variante de a alege mijloace de transport public, prin urmare, dacă utilizatorii ar avea informație necesară la momentul oportun, pot astfel lua o decizie optimă în baza opțiunilor de transport valabile.
\par Acest proiect are scopul de a studia modul în care transportul public este organizat și cum s-ar poate de îmbunătățit calitatea serviciilor și accesul la informație pentru potențialii pasageri. Scopul principal este de a oferi companiilor furnizoare de transport public, servicii prin care informația despre vehicule va putea fi transmisă efectiv către toate entitățile client din cadrul sistemului. Pentru moment, în cadrul lucrării de față accentul se pune pe datele de localizare a vehiculelor, dar același mecanism ar putea fi extins pentru transmiterea mai multor date pe viitor. Din cauza faptului că în calitate de clienți pot fi însăși pasagerii / utilizatorii, dar și companiile de transport, va fi nevoie de luat decizii arhitecturale care permit integrarea mai multor tip de entități client și scalabilitatea sistemului.
